% Transcription — fonds Alexandre Grothendieck, shelfmark 121, batch 3 (pages 41-60).
% Established from the facsimile archives/batches/121/batch-03.pdf (a copy of
% 121.pdf fetched through the project's relay,
% https://grothendieck-relay.onrender.com/source/121.pdf; PDF page k of the
% batch is the archivists' page 40+k, no cover sheet), per the skill
% transcribe-grothendieck. The folder is seventy-two pages in four batches
% (1-20, 21-40, 41-60, 61-72), transcribed in parallel; this is the third.
%
% Pass: Opus 5.5 (claude-opus-5-5), 2026-09-26 — first pass, unchecked against the pages by a human.
%
% Numbering. \page{} follows the PDF count (the Montpellier footer), as
% batch 1 does. The archivists' pencil numbers, bottom left, run two ahead
% throughout this batch: page 41 is pencilled 43, page 48 is 50, page 53 is
% 55, page 56 is 58, page 60 is 62. He does not paginate these leaves.
%
% Dating. \dating{1974}, the folder's own date, alone (no group sentence:
% the group range is added only for undated folders). Nothing on these
% twenty pages carries a date.
%
% Pages skipped: 52, a blank cream sheet. No letter appears in this batch:
% the 1974 letter of the inventory is not among pages 41-60.
% Page 47 is a pink cover inscribed « Équisingularité / Cellules » in ink;
% it takes no \page and his words become the section heading, with a \note.
% Pages summarised: none as a whole. Pages 41, 42, 44, 45 and 46 are mostly
% drawings (strata, tubular neighbourhoods, grids of boundary maps): what
% text and formulas they carry is transcribed, the pictures are described
% in one \note each; the grid diagrams of 44 and 46 are not redrawn.
% Pages transcribed: 41-46, 48-51, 53-60. Several stocks of paper and inks:
% squared paper (41), plain sheets in blue or black ink (42-51, 55-58),
% perforated ruled sheets (53-54), pencil (59-60 and parts of 58).
%
% What the batch is about. Four separate runs, none continuing batch 1's
% numbered paragraphs:
% (a) pp. 41-46: drawings and notes on stratified spaces and their
% reconstruction by gluing — strata V_0, V_1, V_2 with boundaries, a tower
% U ⊂ V ⊂ W of regular strata with « équisingularité », the data 1°)-7°)
% (regular spaces, fibrations with regular fibres, manifolds with boundary,
% an isomorphism of boundaries) and the construction V = V° ⊔_{U'} U,
% W = W° ⊔_{V'} V; grids of boundary maps ∂_j V_i° → V_j°; « PL variétés
% à bord fibrées ».
% (b) pp. 47-51, under his cover « Équisingularité / Cellules »:
% « Contextes » L.P.M., SA.P.M., Alg_C, An_C; existence of the coarsest
% strict admissible partition making a locally finite family of closed
% subspaces unions of leaves (induction on dim X, sketched and broken off);
% the order on leaves and its restriction to an open U; filtrations
% X = X^0 ⊃ X^1 ⊃ ... (III); conditions a) b) c) for an admissible
% partition (connected leaves, non-singular leaves, equisingularity) and
% their behaviour on opens; open leaves = maximal leaves; then p. 51
% « Cellules »: admissible partitions (F_i) (conditions I (0)-(iii)) and the
% equivalent description by the closed leaves Phi_i (conditions II (0)-(iv)),
% with the proof of sufficiency.
% (c) pp. 53-54: « Catégorie modérée »: a category M with a functor
% X |-> |X| to compact spaces, axioms M0)-M4) (finite limits, conservative,
% images of tame maps tame, the unit interval with 1, (x+y)/2 and xy
% tame), sub-objects, the factorisation (M3'), a lemma on effective
% epimorphisms; « ind-objets spéciaux » and pseudo-tame locally compact
% spaces; R = "lim" I_n.
% (d) pp. 55-60: the classes M_n of tame compact parts of I^n (I = [0,1]),
% axioms M'0)-M'3), the category M_0 of pairs (n,K) with tame graphs as
% maps, identities and composition, final object, fibre products and
% kernels, M'4) with 1°)-3°); then in pencil (A)-(D) and 5°)-6°): a
% category K with finite limits and a « foncteur exact à gauche,
% conservatif » phi: M -> K, the equivalent data I, M_n with 1°)-3°),
% (3° bis), 4°) a)-b) on finite sups and e ⊔ e -> I_0, amalgamated sums
% along monomorphisms and the search for a tame structure on them,
% injectivity of I_0 and the cartesian description of sub-objects.
% The labels are his and differ from folder 120's (M1)-(M11): M0)-M4) on
% pp. 53-54 and M'0)-M'4) on pp. 55-56 are his own series here, set as he
% writes them (without the opening parenthesis), not regularised.
%
% Notation fixed for the batch: \mathcal{M}, \mathcal{M}_n, \mathcal{M}_0,
% \mathcal{M}_* for his script M; \widetilde{\mathcal{M}} for the ind-
% objects (p. 54); \mathbb{R}, \mathbb{Q}, \mathbb{N}, \mathbb{Z} as he
% outlines them; I, I_n, I_1, I_0 as written; \Phi_i for his closed
% leaves (p. 51) and \Phi_0 ... for the strata of p. 42; \dot{T}_\alpha as
% on p. 48; \mathcal{K} for the category of pp. 58-60 (his « K » there is
% plain and could be a roman K, noted once). |X| for his « | | » functor.
%
% Readings to check first: the drawings of pp. 41, 42, 44-46 and their
% labels; the interlinear additions of p. 48 (after « dim X ») and the
% struck lines at its foot; the left margin of p. 49 (III); the long
% marginal of p. 51; the pencil of pp. 59-60, faint throughout (the left
% margin of p. 59 is left \ill{}); « kif kif » and « l'on » on p. 54.
% Apparatus: counted in the report (\ill{} and \uncertain{} from the
% source, including those inside \struck{} and \note{}).

\documentclass[11pt,a4paper]{article}
% Le préambule commun à toutes les transcriptions du fonds Grothendieck.
%
% Il tient en une idée : **l'incertitude se compose, elle ne se résout pas.**
% Une transcription qui lisse un mot douteux a détruit la seule chose pour
% laquelle on la faisait — un lecteur doit pouvoir savoir, à chaque endroit,
% ce qui était sur le feuillet et ce qui a été ajouté. Les six macros qui
% suivent sont donc l'appareil critique, et non de la décoration.
%
% Les mêmes commandes sont reconnues par `scripts/render.mjs`, qui produit la
% vue de lecture du site. Une macro ajoutée ici doit l'être aussi là-bas,
% faute de quoi le rendu échouera bruyamment — ce qui est voulu : un rendu
% silencieusement faux est pire qu'un rendu absent.

\ProvidesPackage{grothendieck}[2026/08/08 Transcriptions du fonds Grothendieck]

\usepackage{amsmath,amssymb,amsthm}
\usepackage{mathrsfs}
\usepackage{stmaryrd}
% Les diagrammes commutatifs sont partout dans ce fonds : c'est la langue
% même de la géométrie algébrique de Grothendieck, et une transcription qui
% les rendrait en prose ne transcrirait rien.
\usepackage{tikz-cd}
% Français par défaut — c'est la langue des feuillets — mais l'anglais est
% chargé aussi : la traduction et le résumé sont des documents anglais, et la
% typographie française (espaces avant les deux-points) y serait une faute.
% Ces documents basculent par \selectlanguage{english} en tête de corps.
\usepackage[english,french]{babel}
\usepackage{fontspec}
\usepackage{geometry}
\geometry{a4paper,margin=25mm,marginparwidth=28mm}
\usepackage{marginnote}
\usepackage{xcolor}
% ulem, not soul. `soul`'s \st and \ul are built on a character-by-character
% scan that XeLaTeX's font machinery breaks: the document fails with an
% undefined control sequence rather than degrading. ulem does the same two jobs
% and is Unicode-engine safe. `normalem` stops it hijacking \emph, which the
% transcriptions use for Grothendieck's own emphasis.
\usepackage[normalem]{ulem}
\usepackage{hyperref}
\hypersetup{colorlinks=true,linkcolor=black,urlcolor=black,pdfborder={0 0 0}}

% --- Métadonnées du lot -----------------------------------------------------
% Reprises telles quelles par le script de rendu, qui les affiche en tête de
% la vue de lecture. Elles fixent ce que le fichier prétend couvrir : sans
% elles, une transcription flotte sans point d'attache dans un fonds de seize
% mille feuillets.

\newcommand{\folder}[1]{\gdef\grothFolder{#1}}
\newcommand{\batch}[1]{\gdef\grothBatch{#1}}
\newcommand{\leaves}[2]{\gdef\grothFirst{#1}\gdef\grothLast{#2}}
\newcommand{\foldertitle}[1]{\gdef\grothFoldertitle{#1}}
\gdef\grothFolder{?}\gdef\grothBatch{?}\gdef\grothFirst{?}\gdef\grothLast{?}\gdef\grothFoldertitle{}

% --- Le résumé --------------------------------------------------------------
%
% Il ouvre la lecture modernisée, sous le titre « Résumé », et il est composé
% en petit. Ce n'est pas un ornement : c'est le seul passage du document qui
% ne soit pas des mathématiques, et un lecteur doit pouvoir le reconnaître
% avant de l'avoir lu — puis le sauter s'il connaît déjà le sujet, ou s'y
% arrêter s'il ne le connaît pas. Le filet qui le suit marque la reprise du
% corps.

\newenvironment{resume}
  {\small\setlength{\parskip}{0.45em}}
  {\par\medskip\hrule\medskip}

% --- Mots-clés ---------------------------------------------------------------
%
% En anglais, en clôture du résumé de la lecture modernisée : le vocabulaire
% moderne sous lequel chercher ce que les feuillets construisent. Le manifeste
% du site (`npm run manifest`) les extrait pour étiqueter la cote entière —
% cette ligne est la source unique des tags, il n'y a pas de fichier à côté.
\newcommand{\keywords}[1]{\par\smallskip\noindent
  {\footnotesize\textsc{Keywords} --- \textit{#1}}\par}

% --- L'appareil critique ----------------------------------------------------

% Le numéro de feuillet, en marge. C'est l'ancre qui permet de régler un
% désaccord sur une formule en regardant *un* feuillet plutôt qu'une chemise
% de six cents. Le site s'en sert aussi pour tourner les pages du fac-similé
% quand on fait défiler la transcription.
\newcommand{\leaf}[1]{%
  \marginnote{\footnotesize\color{gray!70}\textsf{#1}}%
  \label{leaf:#1}\ignorespaces}

% Illisible : on ne devine pas. Un mot inventé qui se lit comme les autres est
% le pire résultat possible.
\newcommand{\ill}{\textcolor{red!60!black}{[\,\ldots\,]}}

% Lecture douteuse : le mot est donné, et signalé comme incertain.
\newcommand{\uncertain}[1]{\uline{#1}}

% Ajout de l'éditeur — une lettre manquante, un mot sous-entendu.
\newcommand{\add}[1]{\textcolor{blue!50!black}{[#1]}}

% Biffé par Grothendieck lui-même. Ce qu'il a rayé fait partie du document :
% une rature dit souvent où la pensée a changé de direction.
\newcommand{\struck}[1]{\sout{#1}}

% Note du transcripteur, sur l'état matériel du feuillet.
\newcommand{\note}[1]{\par\smallskip{\small\color{gray!60!black}\textit{[#1]}}\par\smallskip}

% Note marginale de Grothendieck, distincte du corps du texte.
\newcommand{\marginal}[1]{\par\smallskip\noindent\hspace{1em}%
  {\small\color{gray!60!black}#1}\par\smallskip}

% --- Titre ------------------------------------------------------------------

\AtBeginDocument{%
  \begin{center}
    {\large\bfseries Fonds Alexandre Grothendieck --- cote n\textsuperscript{o}~\grothFolder}\\[2pt]
    {\itshape\grothFoldertitle}\\[4pt]
    {\small feuillets \grothFirst--\grothLast\ (lot \grothBatch)}\\[2pt]
    {\footnotesize\color{gray!60!black}Transcription établie d'après le fac-similé numérisé
     par l'Université de Montpellier.\\
     Les lectures incertaines sont soulignées, les illisibles notées [\,\ldots\,],
     les ajouts entre crochets.}
  \end{center}
  \vspace{1em}}

\endinput


\folder{121}
\batch{3}
\pages{41}{60}
\dating{1974}
\watermark{Édition de démonstration}
\foldertitle{Topologie modérée : notes manuscrites (s.d.), lettre (1974)}

\begin{document}

\page{41}
\note{feuillet de papier quadrillé, presque entièrement occupé par des
dessins ; on ne transcrit que les légendes. En haut, un point $V_0$ ; puis
une courbe $V_1$ aboutissant à $V_0$, avec la légende « $\partial V_1 =
V_{1,0}$ » et le schéma $\partial V_1 = V_{10} \to V_0$, $\partial V_1 \to
V_1$. Plus bas, une figure plane : une région hachurée $V_2$, bordée par la
courbe $V_1$, avec des bandes marquées $V_{21}$, $V_{210}$, $V_{20}$ (un
demi-disque autour du point $V_0$). Le reste de la page est blanc.}

\page{42}
\note{page de schémas, en partie au crayon ; on donne ce qui se lit.}

$\Phi_0$ ens fini \qquad $\Phi_1 = \lbrace \Phi_1^{\circ} ; \Phi_1^{b}
\rbrace$ \struck{\ill{}} « 1-spécial »\note{lecture de « spécial »
douteuse ; plusieurs mots biffés à droite, illisibles}

\struck{\ill{} \ill{}} $\mathbb{Z}/2\mathbb{Z}$-torseur (arêtes orientées
\uncertain{ordinaires})

$\Phi_2 = \lbrace F \xrightarrow{i} G \rbrace$ \ill{} « 1-spécial » $\to$
groupoïde ; \uncertain{foncteurs} \ill{} admissibles

$\partial_0 \Phi_2 \in \Pi_0 F$ \qquad $\Psi_0 =$ ens fini, \quad $\Psi_1 =
\lbrace \Psi_1^{\circ} ; \Psi_1^{b} \rbrace$

\begin{tikzcd}[column sep=small, row sep=small]
F_1 & \partial F_1 = F_1|\partial\Phi_1 \arrow[l] \\
\Phi_1 & \partial\Phi_1 \arrow[l, "\pi_1"'] \arrow[d, "\overline{\pi}_1"] \\
 & \Phi_0
\end{tikzcd}

\note{à côté, des schémas analogues pour le niveau 2 : $F_2$, $\partial F_2 =
F_2|\partial\Phi_2$, $\Phi_2 \xleftarrow{\partial_2} \partial\Phi_2$, avec
les composantes $\partial_0\Phi_2$, $\partial_1\Phi_2$,
$(\partial_1\Phi_2)_0$, $(\partial_1\Phi_2)_1$, les flèches $\alpha_0$,
$\alpha_1$, $\beta_0$, $\beta_1$, $\pi'_1$, $(\pi_2)_0$, $(\pi_2)_{1,0}$,
$\pi_{2,1,1}$ vers $\Phi_0$, $\Phi_1$, $\partial\Phi_1$ ; en bas, une tour
$X_0 \to X_1 \to X_2$ figurée par des cases $[\Phi_0]$, $[\Phi_1]$,
$[\Phi_2]$. Ces schémas ne sont pas redessinés.}

\page{43}
\[
  U \subset V \subset W
\]
\begin{itemize}
  \item[] $U$ \struck{\ill{}} régulier
  \item[] $V - U$ régulier
  \item[] $W - V$ régulier
  \item[] $(V, W)$ équisingulier le long de $U$\,?
  \item[] $W$ équisingulier le long de $V - U$
\end{itemize}
dim : \quad $r$ \quad $r+s$ \quad $r+s+t$ \qquad reconstruction ;

\struck{Données}
\begin{itemize}
  \item[1°)] $U$ espace régulier (dim $r$)
  \item[2°)] $U'$ fibré sur $U$, \uncertain{fibres} régulières (de dim
  $s-1$), donc $U'$ régulier
  \item[3°)] $V^{\circ}$ variété à bord, de bord $U'$ (dim $r+s$)
  \item[4°)] $V^{\circ\prime}$ fibré sur $V^{\circ}$, à fibres régulières (de
  dim $t-1$) (sa dim est $r+s+t-1$), son bord est $V^{\circ\prime}|\partial
  V^{\circ} = V^{\circ\prime}|U'$ (de dim $r+s+t-2$), fibré sur $U$ à
  fibres régulières
  \item[5°)] $U''$ autre fibré sur $U$ (dim $r+s+t-1$) à fibres des
  \struck{espaces réguliers} variétés à bord (de dim $s+t-1$), donc
  $\partial U''$ = fibré des bords (de dim $r+s+t-2$), fibré sur $U$ de dim
  $s+t-2$\note{dimensions lues sur un bloc serré ; les exposants sont
  incertains}
  \item[6°)] Un isom.\ de $U$-fibrés en fibres régulières $\partial U''
  \simeq \partial V^{\circ\prime}$. On en conclut, par recollement, un
  \struck{\ill{}} espace régulier $V' = V^{\circ\prime}
  \sqcup_{\partial} U''$
  \item[7°)] \struck{espace} variété à bord $W^{\circ}$, de bord $V'$.
\end{itemize}
On construit
\begin{itemize}
  \item[a)] $V = V^{\circ} \sqcup_{U'} U$ et $U \hookrightarrow V$, $V' \to
  V$ ;
  \item[b)] \struck{b)} $W = W^{\circ} \sqcup_{V'} V$ et $V \to W$.
\end{itemize}

\begin{tikzcd}[column sep=small, row sep=small]
V' = V^{\circ\prime} \sqcup_{\partial} U'' \arrow[r] & W^{\circ} \arrow[r]
& W \\
 & V = V^{\circ} \sqcup_{U'} U \arrow[ur] & U \arrow[l, hook']
\end{tikzcd}

\note{le diagramme, à gauche de la page, est redessiné d'après une
esquisse ; les flèches $V' \to V$ et les inclusions des bords n'y sont pas
toutes lisibles}

\page{44}
\note{page de dessins : une suite de figures numérotées 1°) à 6°) (un point
$U_0$, puis $U_0$ et $U_0'$, des « haltères », des régions hachurées
$V^{\circ\prime}$, $U''$), marquées I et II. En haut, la tour $U_0 \subset
U_1 \subset U_2 \subset U_3 \ldots U_n$, avec $U_1^{\circ}$, $U_2^{\circ}$,
$U_3^{\circ}$, $\partial U_1^{\circ} = U_0'$, $\partial U_2^{\circ}$,
$U_1^{\circ\prime}$, $U_0''$, $\partial U_1^{\circ\prime} = \partial U''$,
$U_2^{\circ\prime}$, $\partial U_2^{\circ\prime}$. À droite, un schéma
vertical :}
\[
  U_n \to U_n^{\circ} \leftarrow \partial U_n^{\circ} ; \qquad
  \partial_{n-1} U_n^{\circ},\ \partial_{n-2} U_n^{\circ},\ \ldots,\
  \partial_0 U_n^{\circ} \longrightarrow U_{n-1}^{\circ},\
  U_{n-2}^{\circ},\ \ldots,\ U_0^{\circ} = U_0 .
\]

\page{45}
PL \quad variétés à bord fibrées sur v.\ à b.\ à fibres \ill{}

\note{le reste de la page porte, au crayon très pâle, des esquisses du même
genre que la page 44 (haltères, bandes), sans texte lisible}

\page{46}
\[
  \partial_j V_i^{\circ} \to V_j^{\circ} ; \qquad
  \partial(\partial_j V_i^{\circ}) = \bigcup_{k} \partial_j
  V_i^{\circ}|\partial_k V_j^{\circ} \ \cup \ \cdots
\]
\note{l'indice de la réunion et le dernier terme, rendu ici par des points, sont illisibles (\ill{}). La page
porte ensuite deux grands tableaux de flèches : le premier part de
$V_i^{\circ}$ et de $\partial V_i^{\circ}$ et descend en dimension ($n_i$,
$n_i - 1$, \ldots) ; le second est une suite $V_i^{\circ} \leftarrow
V_{ij}^{\circ} \leftarrow V_{ijk}^{\circ} \leftarrow V_{ijkl}^{\circ}
\leftarrow \cdots \leftarrow V_{ijkl\ldots p}^{\circ}$, avec des croix
marquant des carrés cartésiens. Ils ne sont pas redessinés.}

\section*{Équisingularité. Cellules}
\note{titre pris sur le feuillet de garde rose (page 47), où il est écrit à
l'encre, souligné, sur deux lignes : « Équisingularité » puis
« Cellules ». Les mots sont de sa main.}

\page{48}
Contextes
\[
  \left\lbrace
  \begin{array}{l}
  \mathrm{L.P.M.} \\
  \mathrm{SA.P.M.} \\
  \mathrm{Alg}_{\mathbb{C}},\ \mathrm{An}_{\mathbb{C}}
  \end{array}\right.
\]

Soient $X$ donné, avec une famille loc.\ finie de sous-espaces fermés
$T_\alpha \subset X$. Alors il existe une partition \add{\emph{stricte}}
admissible \struck{\ill{}} (i.e.\ avec condition de connexité) de $X$, pour
laquelle les $T_\alpha$ soient réunions de feuilles. Parmi celles-ci, il y en
a une plus grossière que toutes autres. Sa formation commute à la
restriction à \uncertain{des} ouverts.
\marginal{NB c'est une \ill{} \ill{} sur $X$}\note{en marge gauche, en face
de l'énoncé, qu'un trait vertical délimite}

La construction de cette dernière montre qu'il \uncertain{suffit} de
\uncertain{prouver} \uncertain{l'existence} loc\supplied{alemen}t, et la
propriété \ill{} \ill{}\ldots\ \ill{} dans $X$ de
\uncertain{dimension} \uncertain{finie} \add{\ill{} \ill{}}
\struck{\ill{}}. Trivial si $\dim X \leqslant 0$ (unique partition
\uncertain{stricte} admissible OK)
\note{les lignes « La construction \ldots\ dans $X$ » sont surchargées
d'une addition interlinéaire (« Dans les \ill{} \ill{} \ldots ») et d'un
passage biffé ; la lecture est très incertaine}

\struck{Considérons les}

Supposons $\dim X > 0$ (ce qui \uncertain{serait} \ill{}) \ill{} \ill{}.
On se débarrasse \struck{des} \ill{} feuilles \ill{} de la partition
\uncertain{cherchée}, \add{\ill{} \ill{}} \struck{conditions} \ill{}
\struck{des $T_\alpha$ \ill{}}\ \uncertain{multiples}
\struck{i.e.\ $r \geqslant 1$ (\uncertain{facteurs}) $T_{\alpha_1} \cap
T_{\alpha_2} \cap \cdots \cap T_{\alpha_r}$ \uncertain{qui} \ill{}}
$X^1$, comme
\[
  X^1 = \bigcup_{\alpha} \dot{T}_\alpha
\]
Soit
\note{le bas du passage est barré de plusieurs traits obliques et d'un
trait horizontal ; la page s'arrête sur « Soit »}

\page{49}
$U$ ouvert $\subset X$, $i, j \in I$, \quad $\overline{F_i} \cap U \neq
\emptyset$ i.e.\ $F_i \cap U \neq \emptyset$ \struck{\ill{}}
\[
  \overline{F_i} \subset \overline{F_j} \ \overset{?}{\Longleftrightarrow}\
  \overline{F_i} \cap U \subset \overline{F_j} \cap U
\]
$\Longrightarrow$ trivial

$\overset{?}{\Longleftarrow}$
\struck{\ill{}}

\struck{\ill{}} Soit $x \in \overline{F_i}$, au voisinage de $x$ on a
$\overline{F_i} \subset \overline{F_j}$, i.e.\ $\bigl(\overline{F_i} \cap
\overline{F_j} = \bigcup_{k \in K} \overline{F_k}\bigr) = \overline{F_i}$ au
voisinage de $x$.

\struck{\ill{}}

$\forall k$ \quad $\overline{F_k} \subset \overline{F_i}$ égalité au
voisinage de $x$ $\Longrightarrow$ égalité partout\,?

Oui, car $F_k$ est \emph{ouvert} dans $\overline{F_i}$ !

Donc si $I_U = \lbrace i \in I \mid \overline{F_i} \cap U \neq \emptyset
\rbrace \subset I$, alors la relation d'ordre de $I_U$ est celle induite par
$U$.\note{sic ; on attend « induite par $I$ »}

\emph{Cor} Si $F_i \cap U \neq \emptyset$, alors $F_i$ ouvert
$\Longleftrightarrow$ $F_i \cap U$ ouvert (i.e.\ $i$ maximal dans $I$
$\Longleftrightarrow$ $i$ maximal dans $I_U$) ; $\Longrightarrow$ trivial ;
$\Longleftarrow$ soit $i$ maximal dans $I_U$, si $j \in I$, $j \geqslant i$,
alors $j \in I_U$, donc $j = i$, OK.\note{la lecture du corollaire est
incertaine par endroits ; plusieurs mots en sont surchargés}

(III)\note{le numéro est cerclé}
\[
  X = X^0 \supset X^1 \supset X^2 \supset \cdots \supset X^n \supset
  \cdots
\]
\emph{filtration loc.\ finie par des fermés}\note{encadré}, \ $\forall i$,
\struck{\ill{}} $X^{i-1} - X^i$ \uncertain{ouvert} dans \ill{} \ldots\
$V^i = X^i - X^{i+1}$ \ldots\ ; $\forall$ ouvert $U$, l'ens.\ des
composantes connexes de $U \cap V^i$ est loc.\ fini sur $U$, $\forall i$ ; et
$\forall$ comp.\ conn.\ $Z_\alpha$ de $U \cap V^i$, $\overline{Z_\alpha} \cap
U \cap V^j$ est une partie ouverte de $U \cap V^j$.
\marginal{(c'est en fait une condition locale, on peut se borner à $U \ni
x$ \ill{})}\note{les indices de cette filtration sont peu lisibles ; la
lecture « $V^i = X^i - X^{i+1}$ » est donnée sous réserve}

\page{50}
Conditions pour une partition admissible
\begin{itemize}
  \item[(a)] feuilles connexes
  \item[(b)] feuilles non singulières (contexte d'espaces
  \add{\uncertain{lisses}} L.P.M.\ ou SA.P.M.\ \ill{})
  \item[(c)] équisingularité.
\end{itemize}
NB a) n'est pas stable en s'induisant sur un ouvert de $X$ (tandis que la
condition d'admissibilité l'est, ainsi que b) et c)). Une bonne condition
stable est \add{localement} [les fibres \struck{loc.\ connexes}], et
\begin{itemize}
  \item[(C)] $\forall$ ouvert $U$ de $X$, la famille de toutes les
  composantes connexes de tous les $U \cap F_i$ est une partition admissible
  de $U$,
\end{itemize}
\add{c) $\forall i$, la famille des composantes de $U \cap F_i$ est loc.\
finie sur $F_i$ \ill{} $(F_{i\alpha})$} i.e.\ \struck{\ill{}} si
$F_{i\alpha}$ est une composante connexe de $U \cap F_i$ : $F_{j\beta} \cap
\overline{F_{i\alpha}} \neq \emptyset \Rightarrow F_{j\beta} \subset
\overline{F_{i\alpha}}$ (ou encore, $\overline{F_{i\alpha}} \cap F_j \cap
U$ est une partie ouverte de $F_j \cap U$, ou encore de $F_j$)
\note{l'addition « c) \ldots\ $(F_{i\alpha})$ » est interlinéaire ; sa
place exacte dans la phrase est incertaine}

Soit $(F_i)$ partition admissible. $F_i$ ouvert dans $X$
$\Longleftrightarrow$ $\overline{F_i}$ maximale. $F_i$ ouvert \add{dans $X$,
et} $\overline{F_i} \supsetneq \overline{F_j}$ $\Longrightarrow$ $F_i$
\add{$\neq \emptyset$} ouvert \add{dense} dans $\overline{F_j}$ et
$\overline{F_i}$ rare dans $\overline{F_j}$\note{sic, les indices $i$ et $j$
semblent échangés dans la seconde moitié} ; donc $F_i$ ouvert dans $X$
$\Longrightarrow$ $F_i$ maximal. Soit $F_i$ maximal, alors
\[
  X - F_i = \bigcup_{j \neq i} \overline{F_j} \quad \text{(fermé)}
\]
car $\forall j \neq i$, $\overline{F_j} \cap F_i = \emptyset$.
\[
  F_i \text{ ouvert} \Longleftrightarrow X - F_i \ \bigl(=
  \textstyle\bigcup_{j \neq i} F_j\bigr) \text{ fermé} \Longleftrightarrow
  \forall j \neq i,\ \overline{F_j} \subset X - F_i
\]
i.e.\ $\forall j \neq i$, \struck{\ill{}} $\exists k$ tel que
$\overline{F_k} \supsetneq \overline{F_j}$, $k \neq i$
$\Longleftrightarrow$ $i$ maximal\,?\note{fin de page serrée ; lecture
incertaine}

\page{51}
\subsection*{Cellules}
$X$ espace topologique, $(F_i)_{i \in I}$ une partition de $X$ (les $F_i$
sont appelés « \emph{feuilles} »). On dit que la partition est
\emph{admissible} si (I) :\note{les chiffres romains I et II sont écrits en marge, devant une accolade qui embrasse chaque liste}
\begin{itemize}
  \item[(0)] $(F_i)$ est une partition de $X$
  \item[(i)] $(F_i)$ est loc.\ finie \struck{\ill{}}
  \item[(ii)] $\forall i$, $\overline{F_i}$ est une réunion de
  feuilles, et
  \item[(iii)] $F_i$ ouvert dans $\overline{F_i}$ (i.e.\ $F_i$ loc.\
  fermé dans $X$).
\end{itemize}
Les $\overline{F_i}$ sont appelés « \emph{feuilles fermées} ». Quand on a
$\overline{F_i} = \overline{F_j}$, on a $F_i = F_j$, car si
$\overline{F_j} \subset \overline{F_i}$ i.e.\ $F_j \subset \overline{F_i}$,
comme $\overline{F_i}$ est réunion de feuilles, on a $F_j = F_i$, ou $F_j
\subset \overline{F_i} - F_i$, mais \add{dans ce cas} \struck{\ill{}}
« $\overline{F_j} \subset \overline{F_i} - F_i$ », donc $\overline{F_j}
\neq \overline{F_i}$ puisque $F_i \neq \emptyset$. Donc la donnée de la
famille $(F_i)$ équivaut à celle de la famille des $\Phi_i =
\overline{F_i}$. Elle est alors \uncertain{caractérisée} (II) :
\begin{itemize}
  \item[(0)] les $\Phi_i$ \struck{\ill{}} recouvrent $X$
  \item[(i)] $(\Phi_i)$ est loc.\ finie
  \item[(ii)] les $\Phi_i$ fermés
  \item[(iii)] $\forall i, j$, $\Phi_i \cap \Phi_j$ est une réunion
  de $\Phi_k$
  \item[(iv)] $\forall i, j$ avec $\Phi_j \subsetneq \Phi_i$,
  \struck{\ill{}} $\Phi_j$ est rare dans $\Phi_i$.
\end{itemize}
Ces conditions nécessaires pour que $(\Phi_i)$ soit la famille des feuilles
fermées d'une partition admissible, sont aussi suffisantes. En effet,
$\forall i \in I$, soit
\[
  F_i = \Phi_i - \bigcup_{\Phi_j \subsetneq \Phi_i} \Phi_j
\]
les $F_i$ forment une partition de $X$ ($\forall x \in X$, $\exists$ unique
$i$ tel que $x \in F_i$), $F_i$ est loc.\ fermé (\uncertain{grâce} à (i) et
(ii)), $(F_i)$ est loc.\ finie (grâce à (i)), et $\forall i$,
$\overline{F_i} = \Phi_i$ (grâce à (iv)), prouvons que $\Phi_i$ est réunion
de feuilles $F_j$, i.e.\ \struck{Considérons le cas} que $\forall j$, $F_j
\cap \Phi_i \neq \emptyset \Rightarrow F_j \subset \Phi_i$. Mais $F_j \cap
\Phi_i \neq \emptyset \Longrightarrow$ \struck{\ill{}} $\Phi_j \subset
\Phi_i$ (prendre $x \in F_j \cap \Phi_i$, et utiliser la description
marginale de $j$ tel que $x \in F_j$), d'où $F_j \subset \Phi_i$ OK.
\marginal{$x \in F_i$, savoir $i$ tel que $x \in \Phi_i$ et $\Phi_i$ soit
le plus petit de ceux-ci (existe grâce à (iii))}\note{en marge gauche,
entre accolades, reliée à la parenthèse « $\exists$ unique $i$ tel que » ;
c'est la « description marginale » à laquelle renvoie la fin de la page}

\page{53}
\section*{Catégorie modérée}
\note{titre souligné ; feuillet perforé à réglure fine}

A) Catégorie $\mathcal{M}$ avec foncteur
\[
  X \mapsto |X| = \varphi : \mathcal{M} \longrightarrow (\text{Espaces
  compacts})
\]
\begin{itemize}
  \item[{[M0)}] $|\ |$ foncteur « transportable »]
  \item[M1)] Dans $\mathcal{M}$ les $\varprojlim$ finies existent, et
  $\varphi$ est exact à g.
  \item[M2)] Le foncteur $|\ |$ est \add{conservatif} \struck{\ill{}}
  [$\Longrightarrow$ le foncteur est \add{fidèle} \struck{\ill{}}]
\end{itemize}
\note{les chiffres de M1) et M2) sont repassés}

\emph{Cor 1} Pour que $f : X \to Y$ dans $\mathcal{M}$ soit un mono, il faut
et il suffit que $\varphi(f)$ le soit. L'application, pour $X \in
\mathrm{Ob}\,\mathcal{M}$,
\[
  \mathrm{SsObj}(X) \xrightarrow{\ \varphi\ } \mathrm{SsObj}(\varphi(X))
\]
est injective, la relation d'ordre du premier étant induite par
\add{le second}.

\emph{Déf.} Sous-espace ($\mathcal{M}$-)modéré de $|X|$ (stable par int.\
finies). Structure modérée \add{sur un espace compact $K$}. Sous-espace
modéré d'une structure modérée. Application modérée.

\emph{Cor 2} Une application \struck{\ill{}} $f : |X| \to |Y|$ est modérée
ssi son graphe l'est.
\begin{itemize}
  \item[M3)] L'image d'un espace modéré par une application modérée est
  modérée.
\end{itemize}

\begin{tikzcd}[column sep=small, row sep=small]
X \arrow[r, "f"] & Y & & {|X|} \arrow[r, "{|f|}"] & {|Y|} \\
X' \arrow[u] \arrow[r] & Y' \arrow[u, hook] & & {|X'|} \arrow[u]
\arrow[r] & {|Y'|} \arrow[u, hook]
\end{tikzcd}

\note{deux carrés entre crochets, marqués d'une croix au centre ; la
flèche de gauche porte un indice peu lisible}

donc $f$ se factorise de façon \add{est} unique en
\begin{itemize}
  \item[(M3$'$)] $X \xrightarrow{f'} Y' \xrightarrow{i} Y$ avec $|f'|$
  surjectif et $|i|$ injectif i.e.\ $i$ mono.
\end{itemize}
\marginal{forme équivalente}

Est-ce que $f'$ épi (effectif) i.e.\ $X \times_{Y'} X \rightrightarrows X
\to Y'$ exact\,? C'est épi car $|f'|$ l'est, et $|\ |$ conservatif ; soit
alors $X \xrightarrow{g} Z$ tel que $g\,\mathrm{pr}_1 = g\,\mathrm{pr}_2$,
se factorise-t-il par $Y'$\,? \struck{\ill{}} $|g|$ se factorise en $h :
|Y'| \to |Z|$, il suffit de voir que $h$ \uncertain{est modéré}.

\emph{Lemme} Soient $X \xrightarrow{f'} Y'$ dans $\mathcal{M}$ et $H : |Y'|
\to |Z|$ tel que $H \circ |f'|$ modéré \add{$= |g|$}, $|f'|$ surjectif, alors
$H$ modéré i.e.\ provient de $h : Y' \to |Z|$.\note{sic : « $h : Y' \to
|Z|$ »}

En effet, \struck{voit d'abord} le graphe de $H$ est l'image du graphe de
$|g| = H \circ f'$, donc est modéré.

\struck{\ill{}}

Ind-objet \emph{spécial}\,*\ de $\mathcal{M}$ : a) défini par une
\emph{suite}
\[
  X_1 \to X_2 \to \cdots \to X_n \to \cdots
\]
avec morphismes de transition des monos\ ; b) $\forall i$, $\exists\, j
\geqslant i$ tel que $\forall k \geqslant j$, $|X_i| \subset
\mathrm{Int}_{|X_k|} |X_j|$.
\note{l'astérisque est de sa main, au-dessus de « spécial » ; aucun renvoi
correspondant n'est visible. Dans b), un $X_k$ biffé précède l'indice
$|X_k|$}

\page{54}
Alors $\varinjlim_i |X_i|$ est un espace loc.\ compact dénombrable à
l'$\infty$, et les $|X_i|$ y forment un système fond.\ de compacts. Le
foncteur
\[
  \mathfrak{X} \longmapsto |\mathfrak{X}| = \varinjlim_i |X_i|
\]
de la catégorie $\widetilde{\mathcal{M}}$ \add{des ind-objets spéciaux de
$\mathcal{M}$} vers les espaces loc.\ compacts, est exact à g.\ (NB
$\mathcal{M}$ est stable par $\varprojlim$ finies dans
$\widetilde{\mathcal{M}}$), conservatif, peut-être pas transportable --- ce
qui nous permet d'introduire la notion de structure \emph{pseudo}-modérée sur
un espace loc.\ compact (s'il est compact, cela revient à la structure
modérée), de sorte que la catégorie des ind-objets spéciaux devient
équivalente à celle des espaces loc.\ cpts pseudo-modérés. On définit encore
la notion de \emph{morphisme} \add{pseudo-}\struck{modéré}\,modéré, de
sous-ensembles (fermés) pseudo-modérés d'un tel $X$. On a que si $f : X \to
Y$ est \emph{propre} et pseudo-modéré, $f(X)$ est une partie pseudo-modérée
de $Y$.
\marginal{compact.}\note{en marge gauche, en face de « espace loc.\
compact »}

\emph{Donnée supplémentaire} a) Structure pseudo-modérée sur $\mathbb{R}$
pour laquelle les intervalles $I_n = [-n, +n]$
(\struck{$n \in \mathbb{Z}$} $n \in \mathbb{N}^*$) sont des sous-ens.\
(ps.)\,modérés, la multiplication par \add{$x$} \add{$n \neq 0$} étant un
morphisme (donc isomorphisme) ps.\ modéré de $\mathbb{R}$. On aura donc
\[
  \mathbb{R} = \text{``}\varinjlim\text{''}\, I_n \simeq \varinjlim I_1 ,
\]
les morphismes de transition $I_1(n) \to I_1(n+1)$ étant la multiplication
par $\frac{n}{n+1}$. Donc la donnée a) revient à la donnée :
b) Structure modérée sur $I_1 = [-1, +1]$ pour laquelle les applications $t
\mapsto rt$ ($r \in \mathbb{Q}^*$, \struck{\ill{}} $|r| \leqslant 1$) sont
modérées. On suppose de plus
\begin{itemize}
  \item[M4)] 1°) le point $1$ de $I_1$ (ou de $\mathbb{R}$, c'est kif kif)
  est modéré
  \item[] 2°) l'application somme $\mathbb{R} \times \mathbb{R}
  \xrightarrow{x+y} \mathbb{R}$ (i.e.\ $(x,y) \mapsto \frac{x+y}{2}$
  \struck{\ill{}} $I_1 \times I_1 \to I_1$, c'est kif-kif) est modérée
  \item[] 3°) l'application produit $\mathbb{R} \times \mathbb{R}
  \xrightarrow{xy} \mathbb{R}$ (i.e.\ $(x,y) \mapsto xy : I_1 \times I_1
  \to I_1$, c'est kif-kif) est modérée.
\end{itemize}
De 1°) 2°) on conclut que \struck{\ill{} \ill{}} les pts $x \in \mathbb{Q}
\subset \mathbb{R}$ sont modérés, que les translations $T_x$
\struck{\ill{}} \ill{} sont modérées, que les intervalles de $\mathbb{R}$ à
extrémités modérées sont \ill{}.\note{la dernière ligne est coupée par le
bord du feuillet}

\page{55}
Moyennant 1°, 2° les applications modérées de $X$ (espace compact modéré)
dans $\mathbb{R}$ forment un \uncertain{sous-module} sur le corps
$\mathbb{R}_{\mathrm{mod}}$ des pts modérés de $\mathbb{R}$ (contenant
$\mathbb{Q}$) ; moyennant 1°, 2°, 3°) c'est un \uncertain{sous-anneau}
\uncertain{algèbre}.\note{deux mots superposés en fin de phrase}

L'axiome supplémentaire suivant sera \uncertain{vérifié} \uncertain{par} la
suite
\begin{itemize}
  \item[\emph{M5}] Pour tout $X \in \mathrm{Ob}\,\mathcal{M}$, $\exists$
  \add{$n \geqslant 0$ et} un morphisme modéré $X \to \mathbb{R}^n$
  \emph{injectif} [ou encore $X \to I^n$ \emph{injectif}].
\end{itemize}

Considérons, pour tout $n \geqslant 0$, l'ensemble \add{$\mathcal{M}_n$}
des parties compactes modérées de $I^n$ ($I = [0,1]$). Cet ensemble
$\mathcal{M}_*$ satisfait aux conditions
\begin{itemize}
  \item[M$'$0)] les $K \in \mathcal{M}_n$ sont compacts ($\Longleftrightarrow$
  fermés)
  \item[M$'$1)] Stabilité par intersections finies
  \item[M$'$2)] Stabilité par images inverses par applications
  \struck{\ill{}} simpliciales $I^n \to I^m$ ($\mathcal{M}_*$ est un ens.\
  cosimplicial)
  \item[M$'$3)] Stabilité par images directes (--- simpliciales)
\end{itemize}

\struck{\ill{}} On récupère $\mathcal{M}$ et $X \mapsto |X| : \mathcal{M}
\to (\mathrm{Esp\,cpts})$ à \uncertain{équivalence} près, connaissant
$\mathcal{M}_*$, de la façon suivante :
\begin{itemize}
  \item[a)] Une structure \uncertain{modérée} sur un espace compact $X$ est
  une classe d'équiv.\ d'appl.\ continues injectives $X \xrightarrow{i}
  \mathbb{R}^n$ \uncertain{à} image $K \in \mathcal{M}_n$, deux telles
  $(i, K_n)$, $(i', K'_{n'})$ étant équivalentes \struck{\ill{}} si le
  graphe de $i'i^{-1} : K_n \to K'_{n'}$, soit $\Gamma \subset I^{n+n'}$,
  est \struck{modéré} $\in \mathcal{M}_{n+n'}$,
  \item[b)] Si $X$, $Y$ sont munis de structures modérées par $i : X
  \xrightarrow{\sim} K \subset I^n$, $i' : Y \xrightarrow{\sim} K' \subset
  I^{n'}$, $\varphi : X \to Y$ est modéré ssi $\Gamma_{i'\varphi i^{-1}}
  \subset I^{n+n'}$ est $\in \mathcal{M}_{n+n'}$.
\end{itemize}
\marginal{$\mathcal{M}$-structure sur $X$}\note{en marge gauche, avec une
flèche vers « $X \xrightarrow{i} \mathbb{R}^n$ » ; les deux phrases de la
page qui suivent « de la façon suivante » avaient d'abord « $\mathbb{R}^n$ »
puis passent à $I^n$, comme il l'écrit}

\page{56}
Partant de $\mathcal{M}_*$ satisfaisant M$'$1, M$'$2, M$'$3, soit
$\mathcal{M}_0$ la catégorie \add{suivante} \struck{définie des espaces
compacts munis d'une $\mathcal{M}_*$-structure}
\begin{itemize}
  \item[] \struck{$\mathrm{Ob}\,\mathcal{M}$ = espaces compacts $X$ avec
  $\mathcal{M}_*$-structure (définie \ill{})}
  \item[a)] $\mathrm{Ob}\,\mathcal{M}_0$ = couples $(n, K)$, $n \in
  \mathbb{N}$, $K \in \mathcal{M}_n$
  \item[b)] $\mathrm{Hom}((n,K),(n',K'))$ = \struck{parties} \add{ens.\ des}
  applications continues \add{$f$} (continues) de $K$ dans $K'$ telles que
  $\Gamma_f$ ($\in I^{n+n'}$) soit $\in \mathcal{M}_{n+n'}$
  $\simeq$ ens.\ des \struck{parties} $\Gamma \in \mathcal{M}_{n+n'}$,
  telles que $\Gamma \subset K \times K'$ et $\Gamma \to K$ induit par pr.\
  soit bijectif.
\end{itemize}
\note{le « $\mathcal{M}_*$-structure » biffé et une accolade, à gauche,
relient les deux premières lignes à a)}

\emph{NB} Pour les identités, on a $\Gamma_{\mathrm{id}_K} \subset K \times
K \subset I^n \times I^n \simeq I^{2n}$ s'identifie à \struck{\ill{}}
l'image de $K$ par $\delta_n : I^n \to I^{2n}$, et on applique M$'$3.
\struck{Comme} Stabilité par composition, on a si \struck{\ill{}} $(K,-)
\xrightarrow{f} (K',-') \xrightarrow{g} (K'',-'')$, \ldots
\[
  \Gamma_{fg} = p^{-1}(\Gamma_f) \cap q^{-1}(\Gamma_g),
\]
où $p$, $q$ sont les appl.\ can.\ $I^{n+n'+n''} \to I^{n+n'}$ et
$I^{n+n'+n''} \to I^{n'+n''}$, et on applique M$'$1) et M$'$2).
\note{« $\Gamma_{fg}$ » : sic ; le graphe écrit est dans
$I^{n+n'+n''}$, la projection sur $I^{n+n''}$ n'est pas écrite}

On a un foncteur $(K, n) \mapsto K$, $\mathcal{M}_0 \to$ (Esp.\ cpts). Il
est fidèle, et conservatif (appliquer M$'$2) \ill{}). On le rend
transportable, ce qui donne $\mathcal{M} \simeq$ équiv.\ \ill{}
\ill{}.\note{« appliquer M$'$2 » : le chiffre peut être 3}

$\exists$ objet final : $e = (0,$ \struck{\ill{}} $I^0 = \lbrace e
\rbrace)$ : en effet, si $(K,-) \in \mathcal{M}_0$, $K \to e$ \ill{} comme
graphe $K$ lui-même.

$\exists$ produits fibrés. \struck{Soit} En effet, soit \struck{\ill{}}
$(K,-)$, $(K',-')$, et $K \times K' \subset I^{n+n'}$ \ldots\ défini par
$(L,-)$, mais
\note{bas de page surchargé : plusieurs formules biffées, dont on ne
distingue que $K \times K'$ et $I^{n+n'}$}

\page{57}
\ill{} $L \subset K \times K'$ est modérée, donc \struck{\ill{}} $L \to K$
et $L \to K'$ le sont. En effet, \struck{\uncertain{si} \ill{}} $g$ est $K
\times K' \xrightarrow{\mathrm{pr}} K$ \ldots\ est modérée, car le graphe
de $K \times K' \xrightarrow{\mathrm{pr}} K$ dans \ldots\ est
\struck{$\Delta_{K \times K'} \subset I^{\ldots}$}
\[
  (K \times K' \times K) \cap \mathrm{Im}\bigl(I^{n+n'} \to
  I^{n+n'+n}\bigr), \qquad (x, x') \mapsto (x, x', x)
\]
qui est $\in \mathcal{M}_{2n+n'}$ par M$'$1, 2, 3. Donc $g$ modérée
$\Rightarrow$ $f$, $f'$ modérés. Inversement, on a $\Gamma_g \subset
I^{m+n+n'}$, $= p^{-1}(\Gamma_f) \cap q^{-1}(\Gamma_{f'})$, $p :
I^{m+n+n'} \to I^{m+n}$, $q : I^{m+n+n'} \to I^{m+n'}$, \ldots\ et
appliquer M$'$1, 2.\note{« $q^{-1}(\Gamma_{f'})$ » : l'indice est lu
$g$ sur la page ; le contexte demande $f'$}

\struck{La fonction ($K \to K \times K'$ \ill{})} \ldots\ immédiat avec
produits finis.

On voit de façon analogue que si $f, g : (K,-) \rightrightarrows (K',-')$,
alors $\mathrm{Ker}(f,g)$ existe dans $\mathcal{M}_0$ et $\mathcal{M}_0 \to
(\mathrm{Ens})$ \uncertain{y} commute (car $\mathrm{Ker}(f,g) =
r\bigl(p^{-1}(\Gamma_f) \cap q^{-1}(\Gamma_g) \cap \delta(I^{n+n'})\bigr)$
est $\in \mathcal{M}_{n+n'+n'}$ \ldots), \ $f, g : K \to K' \times K'
\subset I^{n+n'}$ \ldots\ ; où $p, q : I^{n+n'+n'} \to I^{n+n'}$ sont les 2
proj., $\delta : I^{n+n'} \to I^{n+n'+n'}$, $r : I^{n+n'+n'} \to I^n$.
\marginal{$x, x'_1, x'_2$ ; \ $x'_1 = f(x)$, $x'_2 = g(x)$, $x'_1 =
x'_2$}\note{en marge droite, à la hauteur de la formule du noyau ; la
parenthèse est surchargée et sa lecture est incertaine}

Ainsi on a M1) M2), et on a aussi M3) grâce à M$'$3). On a l'objet $(I, 1) =
I$ de $\mathcal{M}_0$ \struck{\ill{}} \ldots\ d'où l'axiome M5 par
construction.\note{M1), M2), M3), M5 renvoient aux axiomes de la page 53 et
de la page 55}

Alors les conditions supplémentaires \ill{} avant M4), \uncertain{ou} M4),
équivalent :
\begin{itemize}
  \item[M$'$4)] \struck{\ill{}} $\forall q \in \mathbb{Q}$, $|q| \leqslant
  1$, l'ens $\lbrace (x,y) \mid y = qx \rbrace \subset I \times I$ est $\in
  \mathcal{M}_2$
  \item[] 1°) \struck{$\forall I$} $\lbrace 1 \rbrace \subset I$ est $\in
  \mathcal{M}_1$
  \item[] 2°) l'ens $\lbrace x, y, z \mid z = \frac{x+y}{2} \rbrace \subset
  I^3$ est $\in \mathcal{M}_3$
  \item[] 3°) l'ens $\lbrace x, y, z \mid z = xy \rbrace \subset I^3$
  \ldots
\end{itemize}
\note{la ligne 3°) se termine par un trait, comme « idem »}

\page{58}
(A)\note{la lettre est cerclée ; cette page et les deux suivantes sont en
partie au crayon}
$\mathcal{K}$ catégorie avec $\varprojlim$ finies ; $\mathcal{M}
\xrightarrow{\varphi} \mathcal{K}$
\begin{itemize}
  \item[] foncteur exact à g.\ (avec $\varprojlim$ finies calculées dans
  $\mathcal{M}$)
  \item[] conservatif (donc fidèle)
  \item[] $I_0$ objet de $\mathcal{M}$ tel que $\forall X \in
  \mathrm{Ob}\,\mathcal{M}$, $\exists n \in \mathbb{N}$ et mono $i : X \to
  I_0^n$
\end{itemize}
\note{les trois lignes sont réunies par une accolade}
\struck{$\varphi$ transportable}
\note{sa lettre pour la catégorie cible est un K simple ; on la note
$\mathcal{K}$ pour la distinguer des parties $K$ des pages précédentes}

À équivalence près (pour $(\mathcal{M}, \varphi)$), pour $\mathcal{K}$
donné, on donne équivalemment la donnée
\begin{itemize}
  \item[a)] d'un $I \in \mathrm{Ob}\,\mathcal{K}$
  \item[b)] $\forall n \in \mathbb{N}$, d'un ens $\mathcal{M}_n$ de
  sous-objets de $I^n$ avec les conditions
  \item[] 1°) $\mathcal{M}_n$ stable par intersections finies
  \item[] 2°) $\mathcal{M}_n$ stable par images inverses d'appl.\ simpl.\
  $I^m \to I^n$
  \item[] 3°) Pour $\sigma : I^n \to I^{n'}$ \struck{\ill{}} simpl., et $K
  \in \mathcal{M}_n$, tel que $\sigma|K$ soit mono, $\sigma(K) \in
  \mathcal{M}_{n'}$.
\end{itemize}
NB Si on \ill{} $\varphi$ transportable, on a correspondance avec
$\mathcal{M}$ défini à $\mathcal{K}$-isom près.

(B) Supposons que tt morphisme $X \xrightarrow{f} Y$ de $\mathcal{K}$ se
factorise en $X \xrightarrow{g} Y' \xrightarrow{i} Y$, $g$ épi effectif, $i$
monomorphisme (donc nécessairement $Y' = \mathrm{Im}\,g$\note{sic ; on
attend $\mathrm{Im}\,f$}, donc le plus petit sous-objet de $Y$ par lequel se
factorise $f$) \add{\uncertain{et que les} épi effectifs \uncertain{soient}
universels} \struck{\ill{}}. Considérons la condition plus forte que 3°
\begin{itemize}
  \item[(3° bis)] Comme 3°, mais sans supposer que $\sigma|K$ soit mono.
\end{itemize}
Cette condition est équivalente à : tout sous-objet \ill{} d'un objet de
$\mathcal{K}$ « à structure modérée », image d'un morphisme « modéré », est
modéré \ill{}.\note{fin de phrase peu lisible}

(C) Si dans $\mathcal{K}$ \add{\uncertain{dans les sous-objets d'un objet}}
il y a des sup finis (« réunion »). Considérons
\begin{itemize}
  \item[4°)] a) $\mathcal{M}_n$ stable par sup finis ($\forall n
  \geqslant 0$)
  \item[] b) Il existe un morphisme modéré de $e \sqcup e$ dans $I$
  [défini par deux $(e_1, e_2) \in \mathcal{M}_1$ tels que $e_1 \simeq e$
  et $e_2 \simeq e$] qui soit un monomorphisme.
\end{itemize}
\struck{Cette condition équivaut}

\note{la marge gauche porte, au crayon, une note verticale presque
entièrement illisible, où l'on distingue « $X \to Y$ épi dans
$\mathcal{M}$ » et « $\varphi(h)$ mono »}

\page{59}
Cela équivaut, en termes des $(\mathcal{M}, \varphi)$, à la condition que
l'on ait
\begin{itemize}
  \item[a)] Dans $\mathcal{M}$ les sommes disjointes finies existent, et
  $\varphi$ y commute
  \item[b)] $\exists$ monomorphisme $e \sqcup e \to I_0$.
\end{itemize}

(D) Supposons que dans $\mathcal{K}$ les sommes amalgamées $X \sqcup_Y Z$
existent si $Y \to X$ et $Y \to Z$ mono. Alors pour un morph.\ quelconque
$Y \xrightarrow{f} X$, $X \sqcup_Y X$ existe, car on peut remplacer $f$ par
$Y' \to X$ (l'inclusion de son image). $Y \to X$ est un épi ssi $X \sqcup_Y
X \xrightarrow{\nabla_{X/Y}} X$ est iso. Donc si dans $\mathcal{M}$
\uncertain{les mêmes} types de sommes amalgamées existent, et $\varphi$ y
commute, \uncertain{alors} $\varphi(f)$ est un épi ssi $f$ \uncertain{est}
épi ; si dans $\mathcal{K}$ épi $\Rightarrow$ épi eff., kif kif dans
$\mathcal{M}$ \ill{}.\note{« $X \xrightarrow{f} X$ » sur la page ; lu
$Y \to X$ d'après la suite}

On considère un diagramme dans $\mathcal{M}$, \add{à morphismes modérés}
i.e.\ d'objets modérés (de $\mathcal{K}$)

\begin{tikzcd}[column sep=small, row sep=small]
Y \arrow[r, "i"] \arrow[d, "j"'] & X \\
Z &
\end{tikzcd}

avec $i$, $j$ mono, on cherche $X \sqcup_Y Z$ dans $\mathcal{M}$ avec
$\varphi$ y commutant. Donc on cherche une structure modérée sur $X \sqcup_Y
Z$ \uncertain{de} $\mathcal{K}$, \struck{et un diagr} telle que $X \to X
\sqcup_Y Z$, $Z \to X \sqcup_Y Z$ soient modérés, et que ça satisfasse la
propriété universelle.
\note{« PB universel » sur la page, pour « propriété » ; un second schéma,
$X \sqcup Z \rightrightarrows$ vers $X \sqcup_Y Z$, est esquissé à droite}

\note{la marge gauche porte une note au crayon, illisible}

\page{60}
On suppose que dans $\mathcal{K}$, pour tt diagramme de mono $X
\xrightarrow{i} Y$, $X \xrightarrow{j} Z$, il existe une somme amalgamée $S =$
\struck{$Y$} $Y \sqcup_X Z$, et que de plus $Y, Z \rightrightarrows S$ sont
des mono et $Y \times_S Z \leftarrow X$ iso. Je voudrais aussi
\struck{\ill{}} (\ill{} d'exactitude et sommes) (compte tenu des hyp.\ de
factorisation faites plus haut sur $\mathcal{K}$) qu'il existe un diagramme
cartésien (de mono), (D) :

\begin{tikzcd}[column sep=small, row sep=small]
X \arrow[r, hook, "i"] \arrow[d, hook', "j"'] & Y \arrow[d, hook] \\
Z \arrow[r, hook] & S'
\end{tikzcd}

(et donc on peut prendre $S = Y \sqcup_{S'} Z$ (sup de sous-objets)). Il
résulte de ceci que tout mono de $\mathcal{K}$ est effectif (peut-être pas
universel\ldots) (les sommes amalgamées sont alors universelles (stables
par chgt de base quelconque)).\note{« $Y \sqcup_{S'} Z$ » : sic, pour le
sup de $Y$ et $Z$ dans $S'$}

Supposons maintenant que $X$, $Y$, $Z$ et $i$, $j$ soient modérés. On
cherche sur $S$ une structure modérée telle que $Y \to S$ et $Z \to S$ soient
modérés (\uncertain{lorsque} ça \uncertain{sera} une somme amalg.\ dans
$\mathcal{M}$). Cela revient \uncertain{(sic)} à ce qu'on puisse trouver un
diagramme cartésien (D) avec $S'$ et $Y, Z \to S'$ modérés.

(La condition d'\struck{\ill{}} \ill{} \ldots) \quad $O_1 : e_0 \to I_0$
dans $\mathcal{M}$ \ (cf.\ 4°)\note{« $O_1 : e_0 \to I_0$ dans
$\mathcal{M}$ » est encadré}
\marginal{d'où $\sigma_n : e \to I_0^n$}
\begin{itemize}
  \item[5°)] $I_0$ injectif dans $\mathcal{M}$, i.e.\ $\forall$ mono $Y
  \hookrightarrow X$ dans $\mathcal{M}$, et tt $f : Y \to I_0$, $f$ se
  prolonge à $X$ (NB OPS $X = I_0^n$)
  \item[6°)] $\forall$ mono $Y \hookrightarrow X$, $\exists\, X
  \xrightarrow{f} I_0^n$ tel que $Y = f^{-1}(\sigma_n)$, i.e.\ diag.\
  cartésien
\end{itemize}

\begin{tikzcd}[column sep=small, row sep=small]
Y \arrow[r] \arrow[d, hook'] & e \arrow[d, "\sigma_n"] \\
X \arrow[r, "f"'] & I_0^n
\end{tikzcd}

(NB OPS $X = I_0^n$)

\end{document}
